% !TEX TS-program = pdflatex
% !TEX encoding = UTF-8 Unicode

% This is a simple template for a LaTeX document using the "article" class.
% See "book", "report", "letter" for other types of document.

\documentclass[11pt]{article} % use larger type; default would be 10pt

\usepackage[utf8]{inputenc} % set input encoding (not needed with XeLaTeX)

%%% Examples of Article customizations
% These packages are optional, depending whether you want the features they provide.
% See the LaTeX Companion or other references for full information.

%%% PAGE DIMENSIONS
\usepackage{geometry} % to change the page dimensions
\geometry{a4paper} % or letterpaper (US) or a5paper or....
\geometry{margin=2.5cm} % for example, change the margins to 2 inches all round
% \geometry{landscape} % set up the page for landscape
%   read geometry.pdf for detailed page layout information
\setlength{\headheight}{13.6pt}

\usepackage{graphicx} % support the \includegraphics command and options

% \usepackage[parfill]{parskip} % Activate to begin paragraphs with an empty line rather than an indent

%%% PACKAGES
\usepackage{booktabs} % for much better looking tables
\usepackage{array} % for better arrays (eg matrices) in maths
\usepackage{paralist} % very flexible & customisable lists (eg. enumerate/itemize, etc.)
\usepackage{verbatim} % adds environment for commenting out blocks of text & for better verbatim
\usepackage{subfig} % make it possible to include more than one captioned figure/table in a single float
\usepackage{setspace}


% These packages are all incorporated in the memoir class to one degree or another...

%%% HEADERS & FOOTERS
\usepackage{fancyhdr} % This should be set AFTER setting up the page geometry
\pagestyle{fancy} % options: empty , plain , fancy
\renewcommand{\headrulewidth}{0pt} % customise the layout...
\lhead{Konstruktionsprojekt II, Sommersemster 2024}\chead{}\rhead{\thepage}
\lfoot{Paul Mieskes}\cfoot{}\rfoot{Matrikel-Nr.: 552864}

%%% SECTION TITLE APPEARANCE
\usepackage{sectsty}
\allsectionsfont{\sffamily\mdseries\upshape} % (See the fntguide.pdf for font help)
% (This matches ConTeXt defaults)

%%% ToC (table of contents) APPEARANCE
\usepackage[nottoc,notlof,notlot]{tocbibind} % Put the bibliography in the ToC
\usepackage[titles,subfigure]{tocloft} % Alter the style of the Table of Contents
\renewcommand{\cftsecfont}{\rmfamily\mdseries\upshape}
\renewcommand{\cftsecpagefont}{\rmfamily\mdseries\upshape} % No bold!
\usepackage{subfiles}

%%% END Article customizations

%%% The "real" document content comes below...



\onehalfspacing
\title{Konstruktions Projekt II \\ \Huge \textbf{Mechanische Hebevorrichtung} }
\author{\textbf{Paul Mieskes} - Matrikel-Nr.: 552864 \\ paul.mieskes@tuhh.de \\ Studiengang: Mechatronik - Tutor: Jorrit Bargmann}
\date{SoSe 2024} % Activate to display a given date or no date (if empty),
         % otherwise the current date is printed 

\begin{document}
\begin{figure}
\centering
\includegraphics[width=0.5\textwidth]{pics/tuhh.png}
\end{figure}

\maketitle

\centering
\textbf{Variante gemäß Matrikelnummer}\\

\begin{tabular}{| l r l l |}
	
	\hline
	Vorletzte Ziffer & Neigungswinkel der Hebevorrichtung: & $\alpha = 74$  & Grad\\
	& Zulässige Masse der Ladung: & $m_{zul} = 220 $ & kg \\
	\hline
	Letzte Ziffer & (Rotor-)Drehzahl des Elektromotors: & $ n = 1000 $ & min$^{-1}$ \\
	& Zu überwindende Höhendifferenz: & $ h = 950 $ & mm\\
	\hline
\end{tabular}
\vspace{1.5cm}


\begin{tabular}{ | l |  m{10cm}|}
 \hline 
 & Datum, Unterschrifft Tutor:in \\
 \hline 1. Sprechstunde & \\ &\\
 \hline 2. Sprechstunde & \\& \\
 \hline 3. Sprechstunde & \\ &\\
 \hline 4. Sprechstunde & \\ &\\
 \hline
 
\end{tabular}
\thispagestyle{empty} % Removes Page number
% Ende von titel blatt



\raggedright % ab hier rechts orrientierter text

\newpage
\pagenumbering{roman}
\renewcommand*\contentsname{Inhaltsverzeichnis}
\tableofcontents

\newpage
\pagenumbering{arabic}
\setcounter{page}{1} % Reset Page number count

%\subfile{sections/dummy.tex}
 
\subfile{sections/formelzeichen.tex}
\subfile{sections/ausgangssituation.tex}
\subfile{sections/berechnung_alg_groesen.tex}


\end{document}
